\documentclass[11pt]{ctexart}
\usepackage[margin=1in]{geometry}
\usepackage{booktabs}
\usepackage{adjustbox}
\usepackage{amsmath}
\usepackage{setspace}
\usepackage{hyperref}
\hypersetup{colorlinks=true, linkcolor=black, urlcolor=blue}
\setstretch{1.18}

\title{实证思路与当前结果备忘}
\author{}
\date{\today}

\begin{document}
\maketitle

\section*{一、研究目标}

本文拟检验“价格型阶级妥协”的经验机制：在工资增长受限的条件下，低价贸易品和全球供应链在一定时期内压低劳动力再生产成本，从而缓冲实际购买力压力；当关税冲击和供应链冲击削弱这一低价缓冲时，价格压力沿投入产出网络传导，名义工资未能同步补偿，实际购买力和价格型阶级妥协的物质基础受到侵蚀。

本文不把“贸易品价格影响国内价格”作为唯一原创贡献，因为已有文献已经检验进口价格、关税和产业网络对国内价格的影响。本文的增量在于：将这些价格机制置于阶级妥协和劳动力再生产成本的理论框架中，并用产业网络冲击检验该机制是否在美国 2016--2025 年数据中成立。

\section*{二、数据来源}

\begin{enumerate}
  \item BEA 投入产出数据：使用 2019 年 BEA IO 表构造固定产业网络，避免网络结构随冲击期内生变化。
  \item BEA 行业价格数据：使用 2016--2025 年 BEA 行业总产出价格指数和中间品投入价格指数，构造行业层面的价格增长率。
  \item BLS QCEW 工资数据：使用 2016--2025 年私营部门行业工资数据，构造平均年薪、平均周薪及其增长率，并与 BEA summary industry 匹配。
  \item Section 301 关税数据：使用 USTR Section 301 清单，清洗 PDF 中的 HTS8 产品代码，并结合 2017 年中国进口额构造 HS8--NAICS--BEA 行业暴露。
  \item PCE 价格指数：使用 FRED PCEPI 构造全国 PCE 实际工资口径。含年份固定效应时，全国年度 deflator 会被年份固定效应吸收。
  \item 贸易品/非贸易品通胀事实：使用 BLS CPI-U 季调序列。贸易品代理为核心商品，非贸易品代理为核心服务。
\end{enumerate}

\section*{三、清洗与变量构造}

行业网络使用 2019 年 BEA IO 表计算投入份额。网络冲击定义为上游行业关税暴露的投入份额加权和：
\[
NetworkShock_{it}=\sum_j s_{ij,2019}TariffShock_{jt}.
\]

强连接网络参照 Acemoglu et al. 的做法，仅保留投入份额超过 5\% 的行业连接，构造强连接网络冲击变量。关税暴露从 USTR 301 清单提取 HTS8 产品，按实施时间和税率生成年度 active rate，再用 2017 年中国进口额作为权重映射到 NAICS 和 BEA summary industry。主变量采用当期加一年滞后的累计冲击，即 \texttt{t + t-1}。

价格变量包括总产出价格增长和中间品投入价格增长。工资变量包括 QCEW 平均年薪增长、行业价格实际工资增长，以及价格-工资缺口 \(\Delta \ln P_{it}-\Delta \ln W_{it}\)。

\section*{四、计量模型选择}

主模型采用行业和年份双向固定效应：
\[
Y_{it}=\beta NetworkShock_{it}+\alpha_i+\lambda_t+\varepsilon_{it}.
\]

其中 \(Y_{it}\) 分别为行业总产出价格增长、中间品投入价格增长、名义工资增长、实际工资增长或价格-工资缺口。行业固定效应控制行业不随时间变化的结构差异，年份固定效应控制宏观通胀、货币政策、总需求和其他共同冲击。标准误按行业聚类。

双向固定效应本身不能自动识别产业网络传导，因此本文把网络传导显式写入解释变量，即用 2019 年 IO 投入份额加权上游关税暴露。识别并不来自 TWFE 本身，而来自“固定网络结构 \(\times\) 外生关税冲击”的构造。

\section*{五、传导路径}

本文拟检验的路径为：Section 301 关税冲击提高部分上游行业的进口和生产成本；成本冲击沿 2019 年投入产出网络向下游行业传导，尤其通过投入份额较大的关键供应链连接传导；传导首先表现为中间品投入价格和总产出价格上升；名义工资没有同步上升，实际工资和价格-工资缺口出现方向一致的压力；因此，低价贸易品和低价供应链对劳动力再生产成本的缓冲机制被削弱，价格型阶级妥协的物质基础受到侵蚀。

\section*{六、当前主要结果}

当前结果建议这样定位：第一，主显著结果是价格传导，尤其是中间品投入价格传导。完整 IO 网络冲击对总产出价格的影响在 5\% 水平显著；对中间品投入价格的影响在 1\% 水平显著。第二，强连接网络结果更贴合理论。使用投入份额超过 5\% 的关键 IO 连接后，总产出价格和中间品投入价格仍显著，说明价格压力主要沿关键供应链连接传导。第三，工资项不作为主显著证据。名义工资没有显著上升，价格-工资缺口扩大和行业价格实际工资下降在 10\% 水平边际显著，适合作为辅助证据。第四，更细的 NAICS 6 位工资面板没有产生 5\% 显著工资结果，因此不建议把“工资显著下降”作为论文主结论。

\section*{七、标准回归表}

下列表格均由 R 的 \texttt{fixest::etable} 直接导出。星号规则为：*** \(p<0.01\)，** \(p<0.05\)，* \(p<0.10\)。

\subsection*{表 1：主机制结果}
\begin{adjustbox}{max width=\textwidth}
\input{../tables/advisor_table1_main_mechanism.tex}
\end{adjustbox}

\subsection*{表 2：强连接网络与价格-工资缺口}
\begin{adjustbox}{max width=\textwidth}
\begingroup
\centering
\begin{tabular}{lccccc}
   \tabularnewline \midrule \midrule
   Dependent Variables:                         & Output price growth & Input price growth & Output price minus wage growth & Industry-price real wage growth & Output price growth\\  
   Model:                                       & (1)                 & (2)                & (3)                            & (4)                             & (5)\\  
   \midrule
   \emph{Variables}\\
   Strong-link network tariff shock, t plus t-1 & 0.0637$^{**}$       & 0.0714$^{***}$     & 0.0626$^{*}$                   & -0.0626$^{*}$                   &   \\   
                                                & (0.0269)            & (0.0252)           & (0.0338)                       & (0.0338)                        &   \\   
   Network tariff shock, t plus t-1             &                     &                    &                                &                                 & -0.0201\\   
                                                &                     &                    &                                &                                 & (0.0547)\\   
   Network shock x 2021-2022                    &                     &                    &                                &                                 & 0.2102$^{*}$\\   
                                                &                     &                    &                                &                                 & (0.1214)\\   
   \midrule
   \emph{Fixed-effects}\\
   industry\_fe                                 & Yes                 & Yes                & Yes                            & Yes                             & Yes\\  
   year                                         & Yes                 & Yes                & Yes                            & Yes                             & Yes\\  
   \midrule
   \emph{Fit statistics}\\
   Observations                                 & 594                 & 594                & 559                            & 559                             & 594\\  
   R$^2$                                        & 0.25429             & 0.50836            & 0.26111                        & 0.26111                         & 0.26564\\  
   Within R$^2$                                 & 0.00123             & 0.00417            & 0.00095                        & 0.00095                         & 0.01643\\  
   \midrule \midrule
   \multicolumn{6}{l}{\emph{Clustered (industry\_fe) standard-errors in parentheses}}\\
   \multicolumn{6}{l}{\emph{Signif. Codes: ***: 0.01, **: 0.05, *: 0.1}}\\
\end{tabular}
\par\endgroup

\end{adjustbox}

\subsection*{表 3：价格结果稳健性}
\begin{adjustbox}{max width=\textwidth}
\begingroup
\centering
\begin{tabular}{lcccc}
   \tabularnewline \midrule \midrule
   Dependent Variable: & \multicolumn{4}{c}{Output price growth}\\
   Model:                                       & (1)                   & (2)           & (3)      & (4)\\  
   \midrule
   \emph{Variables}\\
   Direct tariff shock, t plus t-1              & 0.0060                &               &          &   \\   
                                                & (0.0331)              &               &          &   \\   
   Network tariff shock, t plus t-1             &                       & 0.0643$^{**}$ &          &   \\   
                                                &                       & (0.0277)      &          &   \\   
   Network tariff shock excluding own industry  &                       &               & 0.0895   &   \\   
                                                &                       &               & (0.0553) &   \\   
   Strong-link network tariff shock, t plus t-1 &                       &               &          & 0.0637$^{**}$\\   
                                                &                       &               &          & (0.0269)\\   
   \midrule
   \emph{Fixed-effects}\\
   industry\_fe                                 & Yes                   & Yes           & Yes      & Yes\\  
   year                                         & Yes                   & Yes           & Yes      & Yes\\  
   \midrule
   \emph{Fit statistics}\\
   Observations                                 & 594                   & 594           & 594      & 594\\  
   R$^2$                                        & 0.25341               & 0.25440       & 0.25418  & 0.25429\\  
   Within R$^2$                                 & $4.72\times 10^{-5}$  & 0.00138       & 0.00108  & 0.00123\\  
   \midrule \midrule
   \multicolumn{5}{l}{\emph{Clustered (industry\_fe) standard-errors in parentheses}}\\
   \multicolumn{5}{l}{\emph{Signif. Codes: ***: 0.01, **: 0.05, *: 0.1}}\\
\end{tabular}
\par\endgroup

\end{adjustbox}

\subsection*{附表：NAICS 6 位工资检验}
\begin{adjustbox}{max width=\textwidth}
\begingroup
\centering
\begin{tabular}{lcccc}
   \tabularnewline \midrule \midrule
   Dependent Variable: & \multicolumn{4}{c}{Nominal wage growth}\\
   Model:                                  & (1)                   & (2)      & (3)                  & (4)\\  
   \midrule
   \emph{Variables}\\
   Direct Section 301 exposure, t plus t-1 & -0.0018               & -0.0081  & -0.0036              & 0.0008\\   
                                           & (0.0044)              & (0.0097) & (0.0036)             & (0.0135)\\   
   \midrule
   \emph{Fixed-effects}\\
   naics\_fe                               & Yes                   & Yes      & Yes                  & Yes\\  
   year                                    & Yes                   & Yes      & Yes                  & Yes\\  
   \midrule
   \emph{Fit statistics}\\
   Observations                            & 9,115                 & 3,127    & 9,115                & 3,127\\  
   R$^2$                                   & 0.11060               & 0.13854  & 0.21646              & 0.16501\\  
   Within R$^2$                            & $8.58\times 10^{-6}$  & 0.00020  & $3.2\times 10^{-5}$  & $2.74\times 10^{-6}$\\   
   \midrule \midrule
   \multicolumn{5}{l}{\emph{Clustered (naics\_fe) standard-errors in parentheses}}\\
   \multicolumn{5}{l}{\emph{Signif. Codes: ***: 0.01, **: 0.05, *: 0.1}}\\
\end{tabular}
\par\endgroup

\end{adjustbox}

\section*{八、建议写法}

本文的经验结论可以表述为：

\begin{quote}
Section 301 关税冲击通过产业投入产出网络显著推高行业价格，尤其显著推高中间品投入价格；这种传导在关键供应链连接中同样成立。与之相比，工资并未同步补偿价格冲击，价格-工资缺口和行业价格实际工资的方向性结果表明实际购买力承压。本文据此说明，低价贸易品和全球供应链曾经在工资增长有限的条件下缓冲劳动力再生产成本，而当这一价格缓冲被关税和供应链冲击削弱时，价格型阶级妥协的物质基础发生动摇。
\end{quote}

\section*{九、下一步}

建议将表 1 和表 2 作为主实证结果，将表 3 作为稳健性结果，将 NAICS 工资表放入附录，说明工资端结果较弱，避免过度声称。后续还需要补充文献定位：前人已经检验进口价格、关税和网络传导；本文贡献在于解释这些价格机制对劳动力再生产和价格型阶级妥协的含义。

\end{document}
